\section*{ID10030: Kinetic Energy (Rotational): θ}
\paragraph{Metadata.}
PiID: 3975187445455 | RTEID: RTE:P38948.7174:T1.3437 | UUID: 142ddaa9-5fd2-5d71-8a12-293ea5c05357

\paragraph{Canonical Statement.}
\textbackslash{}textbf\{Truth Table\}: \textbackslash{}begin\{itemize\} \textbackslash{}item \$I\_\{\textbackslash{}text\{gyro\}\} \textbackslash{}in \textbackslash{}text\{Sym\}\textasciicircum{}+(3)\$ : \$\textbackslash{}checkmark\$ (positive definite inertia tensor) \textbackslash{}item \$T\_\{\textbackslash{}text\{gyro\}\} \textbackslash{}geq 0\$ : \$\textbackslash{}checkmark\$ (kinetic energy non-negative) \textbackslash{}item Euler angles singularity at \$\textbackslash{}theta = 0, \textbackslash{}pi\$ : \$\textbackslash{}checkmark\$ (known, use quaternions) \textbackslash{}end\{itemize\}

\paragraph{Canonical Equation.}
\[
\theta = 0, \pi
\]

\paragraph{Dictionary.}
Kinetic Energy (Rotational): θ — θ = 0, π

\paragraph{Encyclopedia.}
Kinetic Energy (Rotational): θ is an algebraic definition, written θ = 0, π. It is an algebraic definition expressing the quantity directly in terms of the variables on its right-hand side. It is built from π, defining θ. Within the Trawin Zero Point registry it serves as one element of the framework's mathematical narrative.
